\cotefile{14}
\scanpage{59}
est également caractérisé par la condition que le diagramme suivant soit commutatif

\diagram{\begin{tikzcd}[ampersand replacement=\&,column sep=large,row sep=large]
H^1(A_Y) \arrow[r,"u^{(1)}"] \arrow[d,"\varphi_Y^{\unsure{(1)}}"',"{\substack{\text{isom mod}\\\text{torsion}}}"] \& H^1(B_X) \\
H^1(Y) \arrow[r,swap,"{a(\gamma^{\deleted{\unsure{r}}\unsure{n+m-1|m+m-1}}(z))}"] \& H^{2n-1}(X)(n-1) \arrow[u,swap,"\psi_X"]
\end{tikzcd}}

(où la cohomologie est à coefficients dans $\mathrm{Z}_l$, $l$ premier à la caractéristique). Enfin, $u$ ne dépend que de la classe d'équivalence \uline{numérique} de $\deleted{\unsure{z|\mathfrak{z}}}Z$.

Démonstration. Considérons l'homomorphisme composé $\psi_X\,a(\gamma(z))\,\varphi_X^{(1)}$. Il provient d'une classe de \deleted{\unsure{équival|équivalence}} correspondances algébriques (indépendante de $l$), car les trois facteurs sont de ce type. En vertu de 1.7. il est de la forme $u^{(1)}$, où $u:B_X\longrightarrow A_Y$ est un hom de V.A. Il reste à prouver la formule
\[
u(\operatorname{cl}(D))=S_Y({}^tZ(D))
\]
lorsque $ {}^tZ(D)$ est défini (le fait que cette formule caractérise $u$ est immédiat, car pour tout $\zeta\in\operatorname{Pic}^{\unsure{0}}(X)$, il y a un représentant $D$ tel que $ {}^tZ(D)$ soit défini). On va
% archive sheet number: 58

